\chapter{结论、归属与核验}
\label{app:verification}

\section{核心推导的依赖}
\begin{longtable}{p{0.25\textwidth}p{0.34\textwidth}p{0.32\textwidth}}
\toprule
结论 & 核心条件 & 使用的工具\\\midrule
不可辨别误差下限 & 有限静态情境；控制者仅有允许观察 & 记录归纳；有限类识别\\
覆盖率上界 & 独立公开分配；已知概率表 & 乘积亲和度；似然比较\\
全组织下界 & 相同局部策略；严格正核 & KL链式分解；事件数据处理\\
操作性主定理 & 固定有限$d$；统一似然比界 & 上下界；KL与Hellinger比较\\
规模反例 & 随机线性观察；恢复完整情境 & 覆盖计算；Fano；二点下界\\
校准与表示 & 固定二元线性关系；独立调用噪声 & 行空间；覆盖规划；多数集中\\
通信保真 & 相同后处理；完整观察计数 & 柯西不等式；似然因子化\\
停止与模型误差 & 完整证据前缀；明确误差界 & 非负鞅；平方根似然比\\
\bottomrule
\end{longtable}

\section{与已有数学的关系}
统计实验比较提供后处理与决策能力的基本语言；序贯实验和主动假设检验提供选取观察的形式化问题；换测度、Hellinger亲和度、Fano界和鞅界提供概率工具。二元线性观察的可辨别性来自有限域线性代数，多数纠错使用经典有界集中。

本书把任务相对的冲突对、全局覆盖规划和局部KL规划放在同一个有限组织问题中，给出明确常数的连接，再用共同偏差、规模反例和通信变换检验其含义。这些推导的正确性与是否构成新的学术结果分开核对；文献范围记录在随附元数据中。

\section{有限核验怎样使用}
\nolinkurl{scripts/check_theory.py}使用精确分数、有限枚举和浮点容差检查不可辨别误差、覆盖规划、主动记录KL、规模反例、通信压缩、计数充分性与停止概率。它还生成可复算的数据与表格。

\nolinkurl{scripts/check_manuscript.py}检查章节、引文、交叉引用、环境配对和练习数量。构建报告记录XeLaTeX结果与相关页面的目视检查。一般定理由正文证明承担；有限检查用于发现公式、实现与参数口径之间的不一致。

\section{报告数学进展}
可用三个问题检查一项后续改动：它是否扩大了允许的模型而保持结论；是否收紧了整个策略类别的界；是否给出了更经济的可执行组织。只增加一个名称或一个实例时，应说明它承担的是解释还是数学推广。
